% Canonical D034 normalization, 2026-08-24.
% The PDF primitives below remove volatile build metadata.
\ifdefined\pdfinfoomitdate\pdfinfoomitdate=1\fi
\ifdefined\pdftrailerid\pdftrailerid{}\fi
\ifdefined\pdfsuppressptexinfo\pdfsuppressptexinfo=15\fi
\documentclass[10pt,a4paper]{article}
\usepackage[T1]{fontenc}
\usepackage[utf8]{inputenc}
\usepackage[french]{babel}
\usepackage{lmodern}
\usepackage{microtype}
\usepackage{amsmath,amssymb,mathtools,array,tabularx}
\usepackage[normalem]{ulem}
\usepackage{geometry}
\geometry{top=13mm,bottom=13mm,left=16mm,right=16mm}
\pagestyle{empty}
\setlength{\parindent}{1.4em}
\setlength{\parskip}{0.20em}
\setlength{\emergencystretch}{3em}
\newcommand{\Z}{\mathbf Z}
\newcommand{\Q}{\mathbf Q}
\newcommand{\R}{\mathbf R}
\newcommand{\C}{\mathbf C}
\newcommand{\Gm}{\mathbf G_m}
\newenvironment{compactitem}{\begin{list}{}{\setlength{\leftmargin}{1.6em}\setlength{\itemsep}{0pt}\setlength{\parsep}{0pt}\setlength{\topsep}{0.15em}}}{\end{list}}
\newenvironment{refitems}{\begin{list}{}{\setlength{\leftmargin}{3.2em}\setlength{\labelwidth}{2.7em}\setlength{\labelsep}{0.5em}\setlength{\itemsep}{0.55em}\setlength{\parsep}{0pt}\setlength{\topsep}{0.2em}}}{\end{list}}
\begin{document}
\fontsize{8.65}{10.15}\selectfont
\sloppy
% BEGIN_SOURCE_PAGE printed=23 physical=2

\begin{flushleft}\textsc{SOCIÉTÉ MATHÉMATIQUE DE FRANCE}\\
2e série. Mémoire n\textsuperscript{o} 2, 1980, p. 23-33
\end{flushleft}

\vspace{1.5em}
\begin{center}
{\bfseries CYCLES DE HODGE ABSOLUS ET PÉRIODES DES\\[0.35em]
INTÉGRALES DES VARIÉTÉS ABÉLIENNES}

\vspace{1.4em}
par

\vspace{1em}
Pierre DELIGNE

\vspace{0.5em}
rédigé par J. L. BRYLINSKI
\end{center}

\vspace{1em}
On montre ici comment une théorie convenable des ``motifs'' permet de retrouver certains des résultats de Shimura sur les périodes des variétés abéliennes de type C-M [5]. Il a paru instructif de commencer par une description informelle des propriétés des motifs et de leurs réalisations : c'est ce qui est fait au paragraphe 1. Le paragraphe 2 interprète certains résultats de Shimura à l'aide de motifs. Le paragraphe 3 définit les cycles de Hodge absolus, et la théorie des motifs qui leur est subordonnée. Le paragraphe 4 donne une majoration du degré de transcendance du corps engendré par les périodes d'une variété abélienne quelconque.

\vspace{0.8em}
\noindent\textbf{\S\ 1. \quad \uline{LE LANGAGE DES MOTIFS}}

Soit $k$ un corps de caractéristique $0$. Il existe des ``motifs sur $k$'' plusieurs définitions, dont on ignore si elles sont équivalentes, et ayant chacune

% END_SOURCE_PAGE printed=23 physical=2

\newpage

% BEGIN_SOURCE_PAGE printed=24 physical=3

leurs vertus. Nous préciserons au paragraphe 3 laquelle nous est utile.

Les motifs sur $k$ forment une catégorie additive, munie d'un produit tensoriel. Chaque variété projective et lisse $X$ sur $k$ définit un motif $h(X)$ sur $k$, la \uline{cohomologie motivique} de $X$, et tout motif sur $k$ se déduit d'un motif de cette forme (par découpage en morceaux et twist à la Tate). Enfin, à diverses théories de cohomologie pour les variétés algébriques correspondent des foncteurs ``réalisation'' pour les motifs.

Chaque variété abélienne $A$ sur $k$ définit un motif $T(A)$ sur $k$, son $H_1$ motivique, et le foncteur $A\mapsto T(A)$ est pleinement fidèle. Plonger la catégorie des variétés abéliennes dans celle des motifs a pour intérêt de pouvoir prendre des $\otimes$. Mais les objets obtenus ne sont pas géométriques : il n'y a pas de ``fonctions'' sur un motif.

Le \uline{rang} d'un motif est la dimension de l'une quelconque de ses réalisations, que nous allons maintenant décrire.

Si $k=\C$, la \uline{réalisation de Betti} d'un motif $M$ est un groupe abélien libre de rang fini $H_B(M)$, qui est muni d'une \uline{structure de Hodge}, c'est-à-dire d'une décomposition :
\[
H_B(M)\otimes_{\Z}\C=\bigoplus_{p+q=n}H^{p,q}\qquad(\text{avec }\overline{H^{p,q}}=H^{q,p}).
\]
Illustrons ceci sur le motif $T(A)$, qui est le $H_1$ d'une variété abélienne $A$. On définit $H_B(T(A))=H_1(A,\Z)$, homologie de $A=A(\C)$ vu comme espace topologique.

Soit $\operatorname{Lie}(A)$ l'algèbre de Lie du groupe analytique $A$, et soit $\int:H_1(A,\Z)\to\operatorname{Lie}(A)$ l'homomorphisme ``intégration sur un 1-cycle''. On a une suite exacte de groupes analytiques complexes :
\[
0\longrightarrow H_1(A,\Z)\xrightarrow{\int}\operatorname{Lie}(A)\xrightarrow{\exp}A\longrightarrow0.
\]
La structure de Hodge de $H_1(A,\C)$ est telle que $H_1(A,\C)=H^{0,-1}\oplus H^{-1,0}$. On obtient $H^{0,-1}$ comme noyau de l'application $H_1(A,\C)\to\operatorname{Lie}(A)$, prolongement $\C$-linéaire de

% END_SOURCE_PAGE printed=24 physical=3

\newpage

% BEGIN_SOURCE_PAGE printed=25 physical=4

$\int:H_1(A,\Z)\to\operatorname{Lie}(A)$.

Plaçons-nous maintenant sur un corps de caractéristique $0$ et donnons-nous une clôture algébrique $\bar k$ de $k$. Pour tout nombre premier $\ell$, la \uline{réalisation $\ell$-adique} d'un motif $M$ est un $\Z_\ell$-module libre de type fini $H_\ell(M)$ (muni d'ailleurs d'une action du groupe de Galois de $\bar k$ sur $k$). Lorsque $k=\C$, on dispose d'un isomorphisme :
\[
H_\ell(M)\simeq H_B(M)\otimes\Z_\ell.
\]
Si $A$ est une variété abélienne sur $k$, si $\bar k$ est une clôture algébrique de $k$, la réalisation $\ell$-adique de $T(A)$ est notée $T_\ell(A)$ ; c'est le $\Z_\ell$-module :
\[
T_\ell(A)=\varprojlim_n A_{\ell^n}(\bar k)
\quad\bigl(\text{où }A_{\ell^n}(\bar k)\text{ est le groupe des points d'ordre }\ell^n\text{ de }A(\bar k)\bigr).
\]
Si $k=\C$, l'isomorphisme $T_\ell(A)\simeq H_B(T(A))\otimes\Z_\ell$ est facile à décrire. Pour tout $n$, la projection de $\gamma\otimes x\in H_B(T(A))\otimes\Z_\ell$ sur $A_{\ell^n}(\C)$ se calcule ainsi ; si $m\in H_B(T(A))$ est tel que $x-m\in\ell^n(H_B(T(A))\otimes\Z_\ell)$, l'exponentielle de l'élément $m/\ell^n\in H_1(A,\Q)$ est le point d'ordre $\ell^n$ voulu.

La \uline{réalisation de De Rham} de $M$ est un espace vectoriel sur $k$ de dimension finie $H_{DR}(M)$, muni d'une filtration décroissante $F^\bullet$ (la \uline{filtration de Hodge}). Lorsque $k=\C$, on a un isomorphisme
\[
H_B(M)\otimes\C\simeq H_{DR}(M).
\]
Par cet isomorphisme, la filtration de Hodge est telle que $F^p=\bigoplus_{p'>p}H^{p',q}$. Dans le cas du motif $T(A)$, rappelons [7] l'existence d'une ``extension universelle'' $E$ de $A$ par un groupe vectoriel, qui s'inscrit dans une suite exacte de groupes algébriques :
\[
0\longrightarrow H^1(A,\mathcal O_A)^*\longrightarrow E\longrightarrow A\longrightarrow0.
\]
On a $H_{DR}(T(A))=\operatorname{Lie}(E)$, et la filtration de Hodge est donnée par :
\[
0=F^1\subset F^0=H^1(A,\mathcal O_A)^*\subset F^{-1}=\operatorname{Lie}(E).
\]
Si $k=\C$, on a $H_1(E,\Z)\xrightarrow{\sim}H_1(A,\Z)$, d'où un morphisme :
\[
H_1(A,\C)\xleftarrow{\sim}H_1(E,\C)\xrightarrow{\int}\operatorname{Lie}(E)=H_{DR}(T(A)).
\]
C'est l'isomorphisme voulu.

% END_SOURCE_PAGE printed=25 physical=4

\newpage

% BEGIN_SOURCE_PAGE printed=26 physical=5

\noindent\uline{Exemples} : \textbf{1)} Ici, $k=\Q$. On va définir le \uline{motif de Tate} $\Z(1)$ ; c'est le $H_1$ du groupe algébrique $\Gm$. Pour trouver ses diverses réalisations, on substitue $\Gm$ à $A$ dans la définition de $T(A)$. On obtient :
\begin{compactitem}
\item[$-$] $H_B(\Z(1))=2\pi i\cdot\Z$, noyau de $\exp:\C\longrightarrow\C^\times$ ;
\item[] $H_B(\Z(1))\otimes\C=H^{-1,-1}$ ;
\item[$-$] $H_\ell(\Z(1))=\Z_\ell(1)=T_\ell(\Gm)=\varprojlim_n\mu_{\ell^n}(\bar\Q)$ ;
\item[$-$] $H_{DR}(\Z(1))=\Q$.
\end{compactitem}
L'isomorphisme $H_B(\Z(1))\otimes\C\xrightarrow{\sim}H_{DR}(\Z(1))\otimes\C$ envoie $(2\pi i)\otimes1$ sur $2\pi i$. On rappelle les isomorphismes $H_1(\Gm)\simeq H_1(\mathbf P^1-\{0,\infty\})\simeq H_2(\mathbf P^1)$. De plus, pour toute courbe algébrique propre lisse et connexe sur $\Q$, $H_2(X)$ est isomorphe à $H_2(\mathbf P^1)$ (donc à $\Z(1)$). Pour tout entier $d$, on définit un motif $\Z(d)$ de sorte que $\Z(d)=\Z(d-1)\otimes\Z(1)$. On a $\Z(m+n)\simeq\Z(m)\otimes\Z(n)$ pour $m,n\in\Z$.

\textbf{2)} Soit $E$ une courbe elliptique sur $\Q$. Le motif $T(E)$ est de rang 2. Donc le motif $\bigwedge^2T(E)$ est de rang 1, et on a
\[
\bigwedge^2T(E)=\bigwedge^2H_1(E)=H_2(E)=\Z(1).
\]
Précisons l'isomorphisme $\bigwedge^2T(E)\simeq\Z(1)$ dans les différentes réalisations.

Sur $H_1(E(\C),\R)\simeq\operatorname{Lie}(E)$, la structure complexe définit une orientation. Si $(e_1,e_2)$ est une $\Z$-base orientée de $H_B(T(E))=H_1(E(\C),\Z)$, alors $e_1\wedge e_2$ est une base de $H_B(\bigwedge^2T(E))$ ; on lui associe $2\pi i\in H_B(\Z(1))$.

L'isomorphisme $\bigwedge^2(T_\ell(E))\simeq\Z_\ell(1)$ est donné par la forme alternée de Weil [6]. Enfin, on a :
\[
\bigwedge^2H_{DR}(T(E))=\operatorname{Hom}(H^2_{DR}(E),\Q).
\]
Soient $\omega$ et $\eta$ les formes différentielles habituelles de première et deuxième espèce. $H^2_{DR}(E)$ est engendré par le cup-produit $\omega\wedge\eta$ de leurs classes de cohomologie. Par l'isomorphisme
\[
H_B(\bigwedge^2T(E))\otimes\C\longrightarrow H_{DR}(\bigwedge^2T(E))\otimes\C,
\]
$e_1\wedge e_2$ a pour image l'homomorphisme qui envoie $\omega\wedge\eta\in H^2_{DR}(E)$ sur
\[
\begin{vmatrix}\omega_1&\omega_2\\\eta_1&\eta_2\end{vmatrix}\in\C.
\]
La compatibilité des réalisations de Betti et de De Rham de l'isomorphisme $\bigwedge^2(T(E))\simeq\Z(1)$ équivaut à la relation de Legendre
\[
\omega_1\cdot\eta_2-\omega_2\cdot\eta_1=2\pi i.
\]

% END_SOURCE_PAGE printed=26 physical=5

\newpage

% BEGIN_SOURCE_PAGE printed=27 physical=6

\textbf{3)} Soit $A$ une variété abélienne sur $\Q$. A la classe de cohomologie d'une section hyperplane de $A$ dans un plongement projectif, correspond un morphisme de motifs $\bigwedge^2T(A)\to\Z(1)$. En cohomologie de Betti, ceci correspond à une forme alternée $2\pi i\cdot E$ sur $H_B(T(A))$, à valeurs dans $2\pi i\cdot\Z$. En cohomologie de De Rham, on a de même une forme alternée $E_{DR}$ sur $H_{DR}(T(A))$ à valeurs dans $\Q$. Il est classique que $F^0$ est un sous-espace totalement isotrope. On en déduit qu'on peut trouver des différentielles de première espèce $\omega_1,\ldots,\omega_g$ sur $A$, des différentielles de deuxième espèce $\eta_1,\ldots,\eta_g$ sur $A$, telles que l'on ait :
\[
E_{DR}(x,y)=\sum_{i=1}^g\bigl((\omega_i,x)(\eta_i,y)-(\omega_i,y)\cdot(\eta_i,x)\bigr).
\]
La compatibilité entre $E$ et $E_{DR}$ se traduit comme suit. Soit $(e_1,\ldots,e_{2g})$ une base de $H_B(T(A))$, soit $\Omega$ la matrice $n\times2n$ telle que $\Omega_{ij}=\int_{e_j}\omega_i$, et $N$ la matrice $n\times2n$ telle que $N_{ij}=\int_{e_j}\eta_i$. On a :
\[
{}^tN\cdot\Omega-{}^t\Omega\cdot N=2i\pi\cdot E.
\]
Enfin, pour la réalisation $\ell$-adique on trouve une forme alternée sur $T_\ell(A)$ à valeurs dans $\Z_\ell(1)$ : c'est la forme alternée de Weil associée au diviseur de la section hyperplane [6].

\vspace{0.5em}
\noindent\textbf{\S\ 2. \quad \uline{MOTIFS LIÉS AUX VARIÉTÉS ABÉLIENNES DE TYPE C-M}}

Soit $k$ un sous-corps algébriquement clos de $\C$ -- par exemple $\bar\Q$ -- et soit $\mathcal A$ la sous-catégorie de la catégorie des motifs sur $k$, engendrée par les $T(A)$ des variétés abéliennes (pour les opérations de produit tensoriel, passage au dual et facteurs directs). Le foncteur ``réalisation de Betti de motif déduit de $M/k$ par extension des scalaires à $\C$'' envoie $\mathcal A$ dans la catégorie des structures de Hodge. Il résulte formellement du Théorème 1 du \S\ 3 que ce foncteur est pleinement fidèle.

Soient $E$ un corps CM, $S$ l'ensemble de ces plongements complexes et $-:S\to S$ la conjugaison complexe. A chaque fonction $\varphi:S\to\Z$ vérifiant

% END_SOURCE_PAGE printed=27 physical=6

\newpage

% BEGIN_SOURCE_PAGE printed=28 physical=7

\[
(*)\qquad \varphi(\sigma)+\varphi(\bar\sigma)\quad\text{est indépendant de }\sigma\in S,
\]
attachons une structure de Hodge $H_\varphi$, munie d'une structure de module sur l'anneau $\mathcal O_E$ des entiers de $E$ :
\[
\begin{aligned}
H_\varphi&=\mathcal O_E,\quad\text{muni de la structure de module évidente, et dans la décomposition}\\
H_\varphi\otimes\C&=E\otimes\C\simeq\C^S,\quad\text{le facteur $\C$ d'indice $\sigma\in S$ est de type de Hodge}\\
&\hspace{8em}(\varphi(\sigma),\varphi(\bar\sigma)).
\end{aligned}
\]
On vérifie que les $H_\varphi$ sont dans l'image de $\mathcal A$, et même dans l'image de la sous-catégorie engendrée par les variétés abéliennes de type CM. Ceci permet de définir le motif $M_\varphi$ sur $\bar\Q$ comme étant celui de réalisation de Betti $H_\varphi$.

La réalisation de De Rham $H_{DR}(M_\varphi)$ est un espace vectoriel sur $\bar\Q$, muni d'une action de $E$. Pour tout $\sigma\in S$, soit $H_{DR}(M_\varphi)_\sigma$ le sous-$\bar\Q$-espace vectoriel sur lequel $x\in E$ agit comme $\sigma^{-1}(x)\in\bar\Q$. Chacun de ces espaces vectoriels est de dimension 1 engendré par un élément $\omega_\sigma$, et $H_{DR}(M_\varphi)$ en est la somme. Via l'isomorphisme
\[
H_B(M_\varphi)\otimes\C\xrightarrow{\sim}H_{DR}(M_\varphi)\otimes\C,
\]
un élément $e$ de $E=H_B(M_\varphi)\otimes\Q$ correspond à une famille $(e_\sigma)$ avec $e_\sigma\in H_{DR}(M_\varphi)_\sigma$. On vérifie immédiatement que le nombre complexe $p(\sigma,\varphi)$ tel que $p(\sigma,\varphi)e_\sigma=\omega_\sigma$ ne dépend, à multiplication près par $\bar\Q^\times$, ni de $e$ ni du choix de $\omega_\sigma$. Ceci justifie la notation.

Entre les $p(\sigma,\varphi)$ et les \uline{invariants de Shimura} $p_E(\sigma,\varphi)$ (voir [5]), on a la relation :
\[
p_E(\sigma,\varphi)=\frac{p(\sigma,\varphi)}{(2\pi i)^{\varphi(\sigma)}}\quad(\text{modulo }\bar\Q^\times),
\]
(qui n'est qu'une traduction des définitions).

Les identités entre les périodes de Shimura, correspondent à des propriétés algébriques des motifs $M_\varphi$. Avant de les résumer en un tableau, introduisons quelques notations relatives à une extension $F|E$ de corps CM. On note $S_F$ (resp. $S_E$) l'ensemble des plongements complexes de $F$ (resp. $E$). Pour $\sigma\in S_F$, on note $\sigma|E\in S_E$ la restriction à $E$ du plongement $\sigma$.

% END_SOURCE_PAGE printed=28 physical=7

\newpage

% BEGIN_SOURCE_PAGE printed=29 physical=8

Pour $\varphi:S_E\to\Z$ satisfaisant $(*)$, on note $\operatorname{Inf}(\varphi)$ la fonction de $S_F$ vers $\Z$ telle que $\operatorname{Inf}(\varphi)(\sigma)=\varphi(\sigma|E)$. Cette fonction satisfait $(*)$. Pour $\varphi:S_F\to\Z$ satisfaisant $(*)$, on note $\operatorname{Res}(\varphi)$ la fonction de $S_E$ vers $\Z$ telle que
\[
\operatorname{Res}(\varphi)(\sigma)=\sum_{\tau|E=\sigma}\varphi(\tau);
\]
elle satisfait $(*)$.

\begin{center}
\scriptsize
\renewcommand{\arraystretch}{1.18}
\begin{tabularx}{\textwidth}{@{}>{\raggedright\arraybackslash}X|>{\raggedright\arraybackslash}X@{}}
\multicolumn{1}{c}{\textbf{Identités entre périodes}} & \multicolumn{1}{c}{\textbf{Relations entre motifs}}\\\hline
1. $p(\sigma,\varphi'+\varphi'')=p(\sigma,\varphi')\cdot p(\sigma,\varphi'')$ & 1. $M_{\varphi'+\varphi''}=M_{\varphi'}\otimes_E M_{\varphi''}$\\[0.4em]
2. $F|E$ extension de corps CM, $\sigma\in S_F$, $\varphi:S_E\to\Z$ satisfaisant $(*)$ :\\
\quad $p(\sigma,\operatorname{Inf}(\varphi))=p(\sigma|E,\varphi)$ & 2. $M_{\operatorname{Inf}(\varphi)}=M_\varphi\otimes_E F$\\[0.4em]
3. $\sigma\in S_E$, $\varphi:S_F\to\Z$ satisfaisant $(*)$ :\\
\quad $p(\sigma,\operatorname{Res}\varphi)=\displaystyle\prod_{\tau|E=\sigma}p(\tau,\varphi)$ & 3. $M_\varphi$ est un $F$-motif de rang 1, donc un $E$-motif de rang $d=[F:E]$\\
& \quad $M_{\operatorname{Res}(\varphi)}=\bigwedge_E^d M_\varphi$\\[0.4em]
4. $p(\sigma,\uline{1})=2\pi i$ & 4. $M_{\uline{1}}=\Z(-1)$ (dual du motif de Tate)
\end{tabularx}
\end{center}

Ces identités entre périodes permettent de retrouver les théorèmes 1 à 4 de l'exposé de Shimura [5]. Notons que dans [5] Shimura relie les périodes $p_K(\sigma,\varphi)$ aux valeurs en un point spécial de formes modulaires de Hilbert, rationnelles sur $\bar\Q$.

\vspace{0.6em}
\noindent\textbf{\S\ 3. \quad \uline{CYCLES DE HODGE ABSOLUS ET MOTIFS}}

Soit $X$ une variété projective et lisse sur $\C$. Un \uline{cycle de Hodge} de codimension $d$ sur $X$ est un élément de $H^{2d}(X,\Q)$ de type $(d,d)$, ou, ce qui revient au même, un élément de
\[
H^{2d}(X,\Q)(d)\mathrel{\mathop=\limits^{\mathrm{def}}}H^{2d}(X,\Q)\otimes\Z(d)
\]
de type $(0,0)$. Remarquons que pour $x\in H^{2d}(X,\Q)(d)$, il est équivalent de dire que $x$ est de type $(0,0)$ ou que $x$ est dans $F^0$. Une classe de cohomologie d'une sous-variété algébrique de $X$ est un cycle de Hodge. Hodge a conjecturé que tout cycle de Hodge serait combinaison

% END_SOURCE_PAGE printed=29 physical=8

\newpage

% BEGIN_SOURCE_PAGE printed=30 physical=9

linéaire de classes de cohomologie de sous-variétés. Cette conjecture paraissant inaccessible, on introduit une notion qui ``arithmétise'' celle de cycle de Hodge.

Les cycles qu'on définit auront tous les attributs visibles des cycles algébriques, à savoir des réalisations en cohomologie $\ell$-adiques et de De Rham. Un \uline{cycle de Hodge absolu} de codimension $d$ sur une variété projective lisse $X$, sur un corps $k$ de caractéristique $0$, sera une collection de classes de cohomologie dans les groupes
\[
H^{2d}(X,\Q_\ell)(d)\quad(\text{resp. }H^{2d}_{DR}(X)(d))
\]
satisfaisant les deux conditions
\begin{compactitem}
\item[$-$] l'élément de $H^{2d}_{DR}(X)(d)$ est dans $F^0$ ;
\item[$-$] pour tout plongement $k\hookrightarrow\C$, les différentes classes de cohomologie correspondent (via les isomorphismes du \S\ 1) à \uline{une} classe rationnelle de la cohomologie de Betti de $X\otimes_k\C$.
\end{compactitem}
On a les inclusions : (cycles algébriques) $\subset$ (cycles de Hodge absolus) $\subset$ (cycles de Hodge). Utilisons cette définition pour préciser la catégorie des motifs sur laquelle on a travaillé dans les \S\ 1 et 2. Soit d'abord $\mathcal V(k)$ la catégorie des variétés projectives et lisses sur $k$. On construit une catégorie \uline{additive} $\mathcal M(k)$, dans laquelle les groupes $\operatorname{Hom}(M,N)$ sont des espaces vectoriels sur $\Q$, qui est munie des données suivantes :

\textbf{a)} un produit tensoriel $\otimes$, associatif, commutatif, distributif par rapport à l'addition des objets ;

\textbf{b)} un foncteur contravariant $H^*$ de $\mathcal V(k)$ dans $\mathcal M(k)$, bijectif sur les objets, compatible aux sommes, transformant produits en produits tensoriels.

Dans $\mathcal M(k)$, on a, pour $X$ et $Y$ connexes :
\[
\operatorname{Hom}(H^*(X),H^*(Y))=Z_{h.a}^{\dim(Y)}(X\times Y),
\]
où on note $Z_{h.a}^d(Z)$ l'espace vectoriel sur $\Q$ des cycles de Hodge absolus de codimension $d$ sur $Z$. La construction de la catégorie de motifs $\mathcal M(k)$ se fait alors tout comme dans [2], les cycles algébriques étant remplacés par les cycles de Hodge absolus. Grâce à cette nouvelle définition, les composantes de Künneth de la diagonale de $X\times X$ étant de Hodge absolues, on a dans $\mathcal M(k)$ une décomposition :
\[
H^*(X)=\bigoplus_i H^i(X).
\]

% END_SOURCE_PAGE printed=30 physical=9

\newpage

% BEGIN_SOURCE_PAGE printed=31 physical=10

Le résultat qui suit a fait l'objet d'exposés à l'I. H. E. S. en 1979.

\vspace{0.5em}
\noindent\uline{\textbf{Théorème 1}} : \uline{Tout cycle de Hodge sur une variété abélienne est de Hodge absolu.}

Signalons le

\noindent\uline{\textbf{Corollaire}} (Borovoï) : \uline{Soit $A$ une variété abélienne sur un corps $k$ de caractéristique 0, $\xi$ un cycle de Hodge sur $A$. Les correspondants $\ell$-adiques $\xi_\ell$ de $\xi$ sont fixés par un sous-groupe ouvert de $\operatorname{Gal}(\bar k|k)$.}

\vspace{1em}
\noindent\textbf{\S\ 4. \quad \uline{PÉRIODES DES INTÉGRALES DES VARIÉTÉS ABÉLIENNES}}

Soit $A$ une variété abélienne complexe, $H_1(A,\C)=H^{0,-1}\oplus H^{-1,0}$ la décomposition de Hodge de $H_B(A)$ (cf. \S\ 1). Considérons $GL(H_1(A))$ comme groupe algébrique sur $\Q$. On a un morphisme de groupes algébriques complexes :
\[
\Gm\xrightarrow{\mu}GL(H_1(A,\C))
\]
tel que $\mu(z)$ agisse par l'homothétie de rapport $z$ sur $H^{-1,0}$, par l'identité sur $H^{0,-1}$. Le \uline{groupe de Hodge} $G$ de $A$ est l'adhérence de Zariski du sous-groupe $\mu(\Gm)$ de $GL(H_1(A))$. En d'autres termes, $G$ est un sous-groupe algébrique sur $\Q$ de $GL(H_1(A))$. Une fonction régulière sur un ouvert de Zariski de $G$, et $\Q$-rationnelle, est nulle sur $G$ si et seulement si elle est nulle sur $\mu(\Gm)$. On vérifie aisément que $G$ contient les homothéties de $H_1(A)$. Voici une seconde description de $G$, très utile : on considère tous les tenseurs rationnels de type $(-N/2,-N/2)$ dans les produits tensoriels $\bigotimes^N H_1(A)$. Alors $G$ est le sous-groupe de $GL(H_1(A))$ formé des transformations linéaires $g$ pour lesquelles il existe un scalaire $\mu$ tel que : $g\cdot t=\mu^m t$, pour tout tenseur $t$ de type $(m,m)$. Enfin, on montre que $G$ est le groupe qui fixe tous les cycles de Hodge absolus de type $(0,0)$ dans les espaces de tenseurs mixtes sur $H_1(A)$ (mixtes = covariants et contravariants).

Supposons maintenant $A$ définie sur un sous-corps $k$ de $\C$. Posons $g=\dim(A)$.

\noindent\uline{\textbf{Théorème 2}} : \uline{Soit $\omega_1,\ldots,\omega_g$ des différentielles de première espèce sur $A$, $\eta_1,\ldots,\eta_g$ des différentielles de deuxième espèce sur $A$, telles que}

% END_SOURCE_PAGE printed=31 physical=10

\newpage

% BEGIN_SOURCE_PAGE printed=32 physical=11

\uline{$(\omega_1,\ldots,\omega_g,\eta_1,\ldots,\eta_g)$ soit une base de $H^1_{DR}(A)$. Soit $K$ le sous-corps de $\C$ engendré par $k$ et par les périodes $\int_\gamma\omega_i$, $\int_\gamma\eta_i$ $(1\le i\le g,\ \gamma\in H_1(A,\Z))$ ; le degré de transcendance de $K$ sur $k$ est majoré par la dimension de $G$.}

\noindent\uline{\textbf{Démonstration}} : $K$ est le corps de rationalité de l'isomorphisme
\[
H_{1,B}(A)\otimes_{\Q}\C\xrightarrow{\sim}H_{1,DR}(A)\otimes_k\C
\]
donné par ``intégration'' sur un cycle. Cet isomorphisme préserve les classes de cohomologie des cycles de Hodge absolus. Soit $P$ la variété algébrique formée des isomorphismes qui ont cette propriété. De la définition de $G$, et du théorème 1 résulte que $P$ est un espace principal homogène sous $G$, défini sur $k$, donc $\dim(P)=\dim(G)$. On conclut par le lemme suivant :

\noindent\uline{\textbf{Lemme}} : \uline{Soit $z$ un point de $\mathbb A^N(\C)$ ; le degré de transcendance de $k(z)$ sur $k$ est égal à la dimension de la $k$-adhérence de Zariski de $\{z\}$ ;}

\begin{center}C'est une tautologie.\hfill C.Q.F.D.\end{center}

\noindent\uline{\textbf{Remarque}} : De l'exemple 3 du \S\ 1, résulte que $2\pi i\in K$. On peut d'ailleurs montrer que le degré de transcendance de $K$ sur $k(2\pi i)$ est majoré par $\dim(G)-1$.

\noindent\uline{\textbf{Corollaire}} : \uline{Avec les notations du Théorème 2, supposons de plus $A$ simple de type C-M. Le degré de transcendance de $K$ sur $k$ est majoré par $g+1$.}

\noindent\uline{\textbf{Preuve}} : Il suffit de montrer : $\dim(G)<g+1$. Or $\mu(\Gm)$ commute à l'action sur $H_1(A,\Q)$ du corps $E$. Il en résulte que $G$ commute aussi à cette action, et que $G$ est un tore algébrique. Par ailleurs $G$ préserve, à un facteur près, la forme alternée associée à une section hyperplane (voir l'exemple 3 du \S\ 1). D'où l'assertion.

Remarquons que, dans ce dernier cas, on dispose de la minoration suivante :
\[
\dim(G)>2+\log_2(g),
\]
due à Ribet [3].

% END_SOURCE_PAGE printed=32 physical=11

\newpage

% BEGIN_SOURCE_PAGE printed=33 physical=12

\noindent\uline{\textbf{RÉFÉRENCES}}

\begin{refitems}
\item[{[1]}] P. Deligne : Valeurs de fonctions L et périodes d'intégrales, in Proceedings of symposia in Pure Mathematics, Volume 33, 1978.
\item[{[2]}] M. Demazure : Exposé au séminaire Bourbaki n\textsuperscript{o} 365, in Springer Lecture Notes, 180 (1971).
\item[{[3]}] K. Ribet : Division fields of abelian varieties with complex multiplications, dans ce volume.
\item[{[4]}] G. Shimura : On the derivatives of modular forms, Duke Math. Journal 44 (1977), 365-387.
\item[{[5]}] G. Shimura : The periods of abelian varieties with complex multiplication and the special values of certain zeta functions, dans ce volume.
\item[{[6]}] A. Weil : Courbes algébriques et variétés abéliennes, Hermann 1971 (réimpression).
\item[{[7]}] A. Weil : Variétés abéliennes, in ``Algèbre et Théorie des Nombres'', Colloque international du C.N.R.S., 24, 1950, p.124-127.
\end{refitems}

\vspace{2em}
\begin{flushright}
\begin{minipage}{0.48\textwidth}
P. Deligne\\
Institut des Hautes Etudes Scientifiques\\
91440 Bures-sur-Yvette\\[0.6em]
J.-L. Brylinski\\
Ecole polytechnique, Centre de Mathématiques\\
91128 Palaiseau cedex
\end{minipage}
\end{flushright}

\vfill

% END_SOURCE_PAGE printed=33 physical=12

\end{document}
